\documentclass[11pt,a4paper]{article}

\usepackage[utf8]{inputenc}
\usepackage[T1]{fontenc}
\usepackage[english]{babel}
\usepackage{xcolor}
\usepackage{amsmath,amssymb,amsthm}
\usepackage{mathtools}
\usepackage{geometry}
\geometry{margin=2.4cm}
\usepackage{booktabs}
\usepackage{tabularx}
\usepackage{enumitem}
\usepackage[most]{tcolorbox}
\usepackage{hyperref}
\hypersetup{
  colorlinks=true,
  linkcolor=blue!60!black,
  citecolor=blue!60!black,
  urlcolor=blue!60!black
}

% theorem environments
\newtheorem{theorem}{Theorem}[section]
\newtheorem{proposition}[theorem]{Proposition}
\newtheorem{corollary}[theorem]{Corollary}
\newtheorem{lemma}[theorem]{Lemma}
\newtheorem{conjecture}[theorem]{Conjecture}
\theoremstyle{definition}
\newtheorem{definition}[theorem]{Definition}
\newtheorem{postulate}{Postulate}
\theoremstyle{remark}
\newtheorem{remark}[theorem]{Remark}
\newtheorem{anchor}[theorem]{Physical anchor}

% shortcuts
\newcommand{\Sn}[1]{S^{#1}}
\newcommand{\Szero}{S^{0}}
\newcommand{\Sone}{S^{1}}
\newcommand{\Stwo}{S^{2}}
\newcommand{\Sthree}{S^{3}}
\newcommand{\Sfour}{S^{4}}
\newcommand{\Sseven}{S^{7}}
\newcommand{\R}{\mathbb{R}}
\newcommand{\C}{\mathbb{C}}
\newcommand{\Z}{\mathbb{Z}}
\newcommand{\Hbb}{\mathbb{H}}
\newcommand{\CP}{\mathbb{C}P}
\newcommand{\SU}{\mathrm{SU}}
\newcommand{\U}{\mathrm{U}}
\newcommand{\lP}{\ell_{P}}
\newcommand{\tP}{t_{P}}
\newcommand{\rH}{\widetilde H}
\newcommand{\Dd}{\mathsf{D}}
\newcommand{\golden}{\varphi}

\title{\bfseries Information as Dependently Reconstructible Distinguishability:\\
Carriers, Invariance Theorems, and Operational Capacity\\
on Compact \texorpdfstring{$\Sthree$}{S3} Geometry}

\author{%
\texorpdfstring{$^{1}$}{1}Andrei Preece\thanks{\texttt{andre.preece@web.de}; ORCID: 0009-0004-5978-795X}
\quad and \quad
\texorpdfstring{$^{1}$}{1}Boris Batenin\thanks{\texttt{batenin@web.de}; ORCID: 0009-0000-7027-8583}\\[0.35em]
\small \texorpdfstring{$^{1}$}{1}Independent Researchers
}

\date{Berlin, 2026}

\begin{document}
\maketitle

\begin{abstract}
Most physical accounts of information use the word before defining the object: information is
said to ``emerge,'' to be ``conserved,'' or to be ``fundamental,'' while the primitive it
quantifies is left implicit.  This paper proposes a definition-first framework.  The primitive is
\emph{distinguishability}, formalized on topological carriers; the minimal carrier of an
elementary distinction is $\Szero$.  We prove three structural theorems.  (i)~\emph{Invariance}:
the distinguishability profile $\Dd_{k}(X)=\operatorname{rank}\rH_{k}(X)$ is preserved by every
structure-preserving re-description (homotopy equivalence); elementary distinctions can be
neither created nor erased invertibly.  (ii)~\emph{Attributed erasure}: any change of the profile
requires an explicit non-invertible operation (deletion, identification, change of category),
whose logical irreversibility is the topological counterpart of Landauer's thermodynamic erasure
cost.  (iii)~\emph{Carrier transport}: the suspension functor shifts the profile by one degree
without loss, $\Dd_{k+1}(\Sigma X)=\Dd_{k}(X)$; along
$\Szero\to\Sone\to\Stwo\to\Sthree$ the elementary distinction becomes the fundamental class of
$\Sthree$, the minimal joint carrier of phase, direction, spin, and Hopf structure.  A physical
layer is then added by two postulates---carrier realization and a Planck resolution
grain---yielding an exact, closed-form distinguishability capacity for a compact domain: the
band-limited harmonic space on $\Sthree_{R}$ at grain $\varepsilon$ has dimension
$N=\tfrac16 M(M+1)(2M+1)$ with $M=\lfloor\sqrt{1+R^{2}/\varepsilon^{2}}\rfloor$, asymptotically
$\tfrac13(R/\varepsilon)^{3}$, in exact agreement with the Weyl law.  Causal integration of this
capacity is bounded below by the global accessibility time $t_{\mathrm{coh}}=\pi R/c$.  The
framework does not replace Shannon's theory; it supplies the alphabet that Shannon's theory
presupposes.  All interpretive statements---dependent designation, non-emergence, the Planck
boundary as the lower limit of reconstruction rather than a beginning---are labelled as
interpretation and separated from the theorems.
\medskip

\noindent\textbf{Keywords:} theory of information; distinguishability; carriers; suspension;
$\Sthree$; Planck grain; Landauer principle; compact topology; spectral capacity.
\end{abstract}

\tableofcontents

\section{Introduction}

\subsection{The problem: information named before it is defined}

The word ``information'' now carries load in thermodynamics, quantum theory, cosmology, and the
foundations of computation.  Yet the standard quantitative theories are measure theories over
\emph{presupposed} alphabets.  Shannon's theory~\cite{Shannon1948} quantifies the statistics of a
source once a set of distinguishable symbols is given; Kolmogorov
complexity~\cite{Kolmogorov1965} quantifies description length once a machine model and a symbol
set are given; quantum information theory quantifies state distinguishability once a Hilbert
space is given.  None of them answers the prior questions: what is a distinction, what does it
take for a distinction to be physically realized, what operations preserve it, what operations
destroy it and at what price, and how much distinguishability a given physical domain can carry
at a given resolution.

Foundational slogans have tried to close the gap from the other side.  Wheeler's ``it from
bit''~\cite{Wheeler1990} makes the bit primitive and matter derivative; Landauer's ``information
is physical''~\cite{Landauer1991} makes the carrier primitive.  Both intuitions are correct in
part, and both are incomplete without a formal object that says what a ``bit'' is before any
physics, and what physics must add before that object becomes operational.

This paper supplies such an object and develops its consequences with an explicit claim-status
discipline.  The primitive is \emph{distinguishability}, carried by topological spaces; the
minimal carrier of one elementary distinction is the zero-sphere $\Szero=\{-1,+1\}$.  The
framework then proceeds in three layers:

\begin{enumerate}[label=(\Roman*)]
\item a \textbf{mathematical layer} (Secs.~\ref{sec:distinguishability}--\ref{sec:carriers}),
in which invariance, attributed-erasure, and carrier-transport theorems are proved;
\item an \textbf{operational layer} (Secs.~\ref{sec:physical}--\ref{sec:capacity}), in which two
postulates---carrier realization and a resolution grain---turn the mathematics into exact
capacity and accessibility statements on compact $\Sthree$ geometry;
\item an \textbf{interpretive layer} (Sec.~\ref{sec:interpretation}), explicitly labelled, in
which the ontological reading (dependent designation, non-emergence, the Planck boundary as a
limit of reconstruction) is stated without being promoted to theorem.
\end{enumerate}

The companion papers of this programme supply the geometric infrastructure: the suspension-chain
analysis of $\Sthree$ as minimal spin/Hopf carrier~\cite{PB2026susp}, the Ricci-type
reconstruction flow with its no-go theorems and Planck-barrier invariant
region~\cite{BP2026ricci}, and the compact causal accessibility structure of light on
$\Sthree_{R}$~\cite{PB2026light}.  The present paper is the information-theoretic keystone of
that line: it states what the other papers' carriers carry.

\subsection{Claim-status convention}

\begin{center}
\small
\begin{tabularx}{\textwidth}{@{}lX@{}}
\toprule
\textbf{Label} & \textbf{Meaning} \\
\midrule
Theorem / Proposition & Proved under stated mathematical assumptions.\\
Physical anchor & Established physics cited for contact, not re-derived here.\\
Postulate & An operational assumption of the physical layer, explicitly adjustable.\\
Ansatz & A proposed quantitative form, not derived from dynamics.\\
Interpretation & Philosophy-of-physics reading; proves no physics.\\
\bottomrule
\end{tabularx}
\end{center}

\subsection{What this paper does not claim}

\begin{tcolorbox}[colback=gray!4,colframe=gray!50,title=Boundary of the paper]
This paper does not replace or amend Shannon's theory, Kolmogorov complexity, or quantum
information theory; it supplies the pre-statistical object those theories presuppose.  It does
not derive physical law from topology; does not derive the value of $\lP$; does not prove that
the observed universe is globally $\Sthree$; does not derive gravitational entropy bounds; and
does not assert any ontological statement as a theorem.  Every formula in the operational layer
is conditional on the stated postulates.
\end{tcolorbox}

\section{The Primitive: Distinguishability Structures}
\label{sec:distinguishability}

\subsection{Elementary distinction}

\begin{definition}[Elementary distinction]
An \emph{elementary distinction} is the two-point space $\Szero=\{-1,+1\}$ with its discrete
topology.  No metric, phase, direction, or causal structure is assumed at this level.  The two
points are not two objects located somewhere; they are the two poles of a single relational
structure, \emph{this / not-this}.  Neither point carries the structure alone: removing either
yields a one-point space, which is no longer a distinction.
\end{definition}

\begin{definition}[Distinguishability profile]
\label{def:profile}
For a nonempty topological space $X$ of the homotopy type of a finite CW complex, the
\emph{distinguishability profile} is
\begin{equation}
\Dd_{k}(X)\;:=\;\operatorname{rank}\rH_{k}(X;\Z),\qquad k\ge0,
\end{equation}
where $\rH_{*}$ denotes reduced singular homology.  $\Dd_{0}(X)$ counts independent elementary
distinctions (one less than the number of path components); $\Dd_{1}$ counts independent cyclic
(phase-type) distinctions; $\Dd_{2}$ directional (enclosure-type) distinctions; and so on.  The
\emph{total content} is $\Dd(X)=\sum_{k\ge0}\Dd_{k}(X)$.
\end{definition}

For the spheres of the suspension chain,
\begin{equation}
\Dd_{k}(\Sn{n})=\delta_{k,n},\qquad \Dd(\Sn{n})=1 :
\label{eq:sphere-profile}
\end{equation}
each sphere carries exactly one unit of distinguishability, concentrated in degree $n$.  In
particular $\Dd_{0}(\Szero)=1$ is the formal trace of the minimal distinction, and
$\Dd_{3}(\Sthree)=1$ is the fundamental class of the compact carrier.

\begin{remark}[Why homology]
The profile could be built from other invariants.  Reduced homology is chosen because it is
(a)~homotopy invariant, so it cannot see mere re-description; (b)~computable; (c)~graded, so it
separates \emph{kinds} of distinction (component / cycle / enclosure) rather than collapsing
them into one number; and (d)~equipped with the suspension isomorphism, which is exactly the
carrier-transport operation the programme needs.  Nothing in the definitions forbids refining
the profile by torsion subgroups or by homotopy groups; rank homology is the minimal choice
sufficient for the theorems below.
\end{remark}

\subsection{Invariance under re-description}

\begin{theorem}[Invariance of distinguishability]
\label{thm:invariance}
If $f:X\to Y$ is a homotopy equivalence, then $\Dd_{k}(X)=\Dd_{k}(Y)$ for all $k$, and $f$
induces a bijection between the path components of $X$ and $Y$.  Consequently no
structure-preserving re-description creates or erases any distinction, elementary or higher.
\end{theorem}

\begin{proof}
A homotopy equivalence induces isomorphisms on all reduced homology groups and a bijection on
$\pi_{0}$~\cite{Hatcher2002}.  Ranks and component counts are therefore equal.
\end{proof}

\begin{corollary}[No invertible erasure, no invertible creation]
\label{cor:no-invertible}
Let $f:X\to Y$ be continuous.  If $f$ merges two path components of $X$ (or if $Y$ has more
path components than $X$ admits under any bijection with $\pi_{0}(X)$), then $f$ is not a
homotopy equivalence: elementary distinctions cannot be removed, and cannot be introduced, by
any invertible operation.
\end{corollary}

\begin{proof}
Merging components makes $\pi_{0}(f)$ non-injective; a component surplus makes it
non-surjective.  Either contradicts the bijectivity required by Theorem~\ref{thm:invariance}.
\end{proof}

\subsection{Attributed erasure: every change of profile has an exhibited operator}

Distinctions can of course be destroyed---but only by operations that leave the
structure-preserving category.

\begin{proposition}[The price of erasing the minimal distinction]
\label{prop:erasure}
The identification map $q:\Szero\to\{*\}$ satisfies $\Dd_{0}:1\mapsto0$ and admits no homotopy
inverse.  Deleting a point likewise destroys the structure ($\Szero$ minus a point is not
$\Szero$), and there is no map at all from the nonempty $\Szero$ to $\varnothing$.  Hence the
minimal distinction can be erased only by an explicit non-invertible operation---identification,
deletion, or exit from the category---never by continuous deformation.
\end{proposition}

\begin{proof}
$\{*\}$ has $\rH_{0}=0$ while $\rH_{0}(\Szero)\simeq\Z$; homotopy equivalences preserve reduced
homology, so $q$ has no homotopy inverse.  The remaining statements are immediate.
\end{proof}

\begin{proposition}[Non-monotonicity: identifications are not pure erasure]
\label{prop:nonmonotone}
Quotient operations need not decrease total content.  Example: the interval $[0,1]$ is
contractible, $\Dd([0,1])=0$, yet the identification $0\sim1$ produces $\Sone$ with
$\Dd_{1}=1$.  Therefore ``erasure'' and ``creation'' are not global monotonicity statements;
they are \emph{per-operation bookkeeping}.  Every change of the profile must be attributed to an
exhibited operator, and the operator's effect must be computed, not assumed.
\end{proposition}

\begin{proof}
$\rH_{1}([0,1])=0$ and $\rH_{1}(\Sone)\simeq\Z$; the quotient map realizes the change.
\end{proof}

\begin{remark}[Interpretation: interdependence cannot be contracted, only failed to be preserved]
Proposition~\ref{prop:erasure} is the formal content of a claim the programme states
ontologically: minimal interdependent distinction is indestructible \emph{within} operations
required to preserve structure.  If a description asserts that a distinction ``has
disappeared,'' the description owes an operator---a quotient, a deletion, a change of
category---and the operator's price is the loss of invertibility.  Nothing in the
structure-preserving language derives the disappearance; it can only be imposed as an act of
forgetting.
\end{remark}

\begin{anchor}[Landauer's principle as the thermodynamic face of attributed erasure]
\label{anchor:landauer}
In physics, the logically irreversible erasure of one elementary distinction (one bit) in an
environment at temperature $T$ dissipates at least $k_{B}T\ln2$ of
heat~\cite{Landauer1961,Bennett1982}.  The topological statement (erasure requires a
non-invertible operator) and the thermodynamic statement (logical irreversibility has a minimum
dissipation) are two faces of one discipline: \emph{nothing erases a distinction for free}.  We
cite this as contact with established physics; we do not claim to derive the constant
$k_{B}T\ln2$ from topology.
\end{anchor}

\section{Carriers and Carrier Transport}
\label{sec:carriers}

\subsection{Carriers}

\begin{definition}[Carrier]
A \emph{carrier} is a space on which a stated package of structures exists naturally and
compatibly.  A carrier is not, by that fact alone, a derivation of dynamics or matter content
(the restricted usage of~\cite{PB2026susp}).
\end{definition}

\subsection{The transport theorem}

\begin{theorem}[Carrier transport; degree-shift conservation]
\label{thm:transport}
The suspension $\Sigma X=(X\times[0,1])/(X\times\{0\},\,X\times\{1\})$ induces natural
isomorphisms
\begin{equation}
\rH_{k}(X)\;\simeq\;\rH_{k+1}(\Sigma X)\qquad\text{for all }k\ge0,
\end{equation}
hence $\Dd_{k+1}(\Sigma X)=\Dd_{k}(X)$ and the total content is preserved:
$\Dd(\Sigma X)=\Dd(X)$.  Suspension changes the carrier and the degree of every distinction;
it neither creates nor destroys any.
\end{theorem}

\begin{proof}
Mayer--Vietoris applied to the two cones covering $\Sigma X$, both contractible, with
intersection homotopy equivalent to $X$~\cite{Hatcher2002}.  Naturality is standard.
\end{proof}

\begin{corollary}[The fundamental class of $\Sthree$ is a suspended distinction]
\label{cor:s3class}
Along $\Szero\xrightarrow{\Sigma}\Sone\xrightarrow{\Sigma}\Stwo\xrightarrow{\Sigma}\Sthree$,
\begin{equation}
\rH_{0}(\Szero)\simeq\rH_{1}(\Sone)\simeq\rH_{2}(\Stwo)\simeq\rH_{3}(\Sthree)\simeq\Z ,
\end{equation}
carrying the distinction class of $\Szero$ to a generator of $H_{3}(\Sthree;\Z)$
\cite{PB2026susp}.  The elementary distinction, the phase cycle, the directional enclosure, and
the compact volume class are one invariant on four carriers.
\end{corollary}

\begin{remark}[Interpretation: carrier changes; generative relation persists]
Theorem~\ref{thm:transport} is the exact form of the programme's central ontological formula:
\emph{information is not produced from non-being; it is transformed through carriers of
distinguishability}.  The invariant is neither simply the same (its degree, its available
structures, its phenomenal unfolding all change) nor simply different (it is the image of one
class under canonical isomorphisms).  A new state of comprehension does not entail a new state
of the invariant content.
\end{remark}

\subsection{Enrichment: what higher carriers add}

Suspension preserves content but changes what the carrier can \emph{do} with it.  On $\Sone$ the
distinction acquires phase ($\U(1)$); on $\Stwo$, direction ($\CP^{1}$, the Bloch sphere); on
$\Sthree$, simultaneously: compactness, simple connectivity, the non-Abelian group $\SU(2)$,
spin structure, parallelizability, and the nontrivial Hopf bundle
$\Sone\hookrightarrow\Sthree\to\Stwo$ with $c_{1}=1$.  The minimality of $\Sthree$ for this
joint package is a theorem of the companion paper~\cite{PB2026susp} and is used here as
infrastructure: $\Sthree$ is the first carrier on which a distinction can be phase-directed,
spin-carried, and globally framed at once, while remaining \emph{loop-trivial but
carrier-nontrivial} ($\pi_{1}=0$, $H_{3}\simeq\Z$).

\begin{remark}[Qubit contact]
Normalized two-level quantum states live on $\Sthree\subset\C^{2}$; physical ray states on
$\CP^{1}\simeq\Stwo$; the projection is the Hopf map, and its nontriviality ($c_{1}=1$) is the
precise sense in which phase and direction do not factorize
globally~\cite{PB2026susp}.  Quantum information theory thus already operates on the minimal
joint carrier identified by the suspension chain.  This is contact, not derivation: probability
rules and measurement maps are additional inputs.
\end{remark}

\section{The Operational Layer: Physical Realization of Distinguishability}
\label{sec:physical}

The mathematical layer says what distinguishability is and how it moves between carriers.  The
physical layer says what it takes for a distinction to be \emph{operational}---measurable,
communicable, thermodynamically accountable.  Two postulates suffice.

\begin{postulate}[Carrier realization]
\label{post:carrier}
A distinction is physically realized only when represented on a carrier equipped with metric and
causal structure.  Without scale there is no distance; without time there is no causal relation;
without causal relation there is no physically unfolded information; without distinguishability
there is no observable phenomenon.  Schematically:
\begin{equation}
\text{distinction}\;\to\;\text{carrier}\;\to\;\text{order/phase}\;\to\;\text{causal relation}
\;\to\;\text{reconstructible information}.
\label{eq:chain}
\end{equation}
\end{postulate}

\begin{postulate}[Resolution grain]
\label{post:grain}
Effective physical reconstruction assigns distinguishable causal-geometric states only at
spatial scales $\ge\lP$ and temporal scales $\ge\tP=\lP/c$.  $\Szero$ is the formal carrier of
minimal distinction; $\lP$ and $\tP$ are its lower metric-causal realization in effective
physics.  Below the grain the model assigns no operationally verifiable state---a boundary of
the reconstruction language, not an ontological wall.
\end{postulate}

\begin{anchor}[Geometric support for Postulate~\ref{post:grain}]
In the companion Ricci-type flow analysis~\cite{BP2026ricci}, a reflecting barrier at
$a=\lP$ defines a forward-invariant region with $|\mathrm{Rm}|^{2}\le C/\lP^{4}$, mirroring a
self-adjoint Wheeler--DeWitt endpoint on $a\in[\lP,\infty)$ with zero probability flux through
$a=\lP$.  The same work proves that natural internal forcings cannot control the singular round
mode, so the grain enters as a \emph{boundary of the effective state space}, not as a smooth
regulator.  Postulate~\ref{post:grain} adopts exactly this reading; it does not claim a
microscopic quantum-gravity derivation of $\lP$.
\end{anchor}

\begin{remark}[Interpretation: probability one, not probability zero]
Under Postulate~\ref{post:grain} the effective physical domain has support on $a\ge\lP$.  The
sub-Planck interval is not another physical branch of probability one; it lies outside the
assigned operational measure.  The question of a ``transition'' from $a<\lP$ into physical
reality does not arise inside the model, because $a<\lP$ is not a state of the model.  Reality
is not reconstructed from probability zero; the effective domain of reconstruction already has
probability-one support on distinguishable causal geometry, bounded below by the Planck grain.
\end{remark}

\section{Capacity Theorems on Compact Carriers}
\label{sec:capacity}

The postulates convert the profile into countable operational alphabets.  On a compact carrier
the count is exact.

\begin{definition}[Operational alphabet at grain $\varepsilon$]
Let $(M,g)$ be a closed Riemannian manifold with Laplace--Beltrami operator $\Delta_{g}$.  The
\emph{band-limited space at grain $\varepsilon$} is
\begin{equation}
\mathcal A_{\varepsilon}(M,g)\;=\;\operatorname{span}\{\,\text{eigenmodes }Y:\;
-\Delta_{g}Y=\lambda Y,\ \sqrt{\lambda}\le1/\varepsilon\,\},
\end{equation}
and the \emph{distinguishability capacity} is
$N_{\varepsilon}(M,g)=\dim\mathcal A_{\varepsilon}(M,g)$: the number of independently
specifiable field degrees of freedom resolvable at scale $\varepsilon$.
\end{definition}

\begin{theorem}[Exact capacity of the round three-sphere]
\label{thm:capacity}
On $\Sthree_{R}$ the scalar Laplacian has eigenvalues $\lambda_{n}=n(n+2)/R^{2}$ with
multiplicities $(n+1)^{2}$, $n=0,1,2,\dots$~\cite{PB2026susp}.  Hence, with
\begin{equation}
M\;=\;\Big\lfloor\sqrt{1+R^{2}/\varepsilon^{2}}\Big\rfloor ,
\end{equation}
the capacity at grain $\varepsilon$ is exactly
\begin{equation}
N_{\varepsilon}(\Sthree_{R})\;=\;\sum_{n=0}^{M-1}(n+1)^{2}
\;=\;\frac{M(M+1)(2M+1)}{6}
\;\sim\;\frac13\left(\frac{R}{\varepsilon}\right)^{3}
\quad(R/\varepsilon\to\infty).
\label{eq:capacity}
\end{equation}
\end{theorem}

\begin{proof}
The condition $\sqrt{\lambda_{n}}\le1/\varepsilon$ reads $n(n+2)\le R^{2}/\varepsilon^{2}$,
i.e.\ $(n+1)^{2}\le1+R^{2}/\varepsilon^{2}$, i.e.\ $n+1\le M$.  Summing the multiplicities gives
$\sum_{j=1}^{M}j^{2}=M(M+1)(2M+1)/6$.  The asymptotic follows from $M\sim R/\varepsilon$.
\end{proof}

\begin{remark}[Weyl consistency check]
The Weyl law for a closed $3$-manifold gives
$N(\lambda)\sim\operatorname{Vol}(M)\,\lambda^{3/2}/(6\pi^{2})$.  With
$\operatorname{Vol}(\Sthree_{R})=2\pi^{2}R^{3}$ and $\lambda=1/\varepsilon^{2}$ this yields
$N\sim\tfrac13(R/\varepsilon)^{3}$, in exact agreement with \eqref{eq:capacity}.  Two elementary
sanity checks: $R/\varepsilon=1$ gives $M=1$, $N=1$ (only the constant mode); $R/\varepsilon=2$
gives $M=2$, $N=5$ (constant plus the four $n=1$ modes).
\end{remark}

\begin{corollary}[Planck capacity of a quasi-local $\Sthree$ domain]
\label{cor:planck-capacity}
Under Postulate~\ref{post:grain} ($\varepsilon=\lP$), a quasi-local domain reconstructed as
$\Sthree_{R}$ carries the finite scalar-field capacity
\begin{equation}
N_{\lP}(\Sthree_{R})\;=\;\frac{M(M+1)(2M+1)}{6},\qquad
M=\Big\lfloor\sqrt{1+R^{2}/\lP^{2}}\Big\rfloor
\;\sim\;\frac{R}{\lP},
\end{equation}
so $N_{\lP}\sim\tfrac13(R/\lP)^{3}$.  For a binary source on this alphabet the classical
capacity is $N_{\lP}$ bits; for an alphabet of $q$ resolvable amplitudes per mode it is
$N_{\lP}\log_{2}q$ bits.  Finiteness is a theorem given the postulates; the identification of
$\varepsilon$ with $\lP$ is the postulate, not the theorem.
\end{corollary}

\begin{remark}[Volume scaling versus area scaling: an honest tension]
\label{rem:holography}
Equation \eqref{eq:capacity} is a \emph{mode count of band-limited fields} and scales with
volume, $(R/\lP)^{3}$.  Gravitational entropy bounds suggest area scaling,
$(R/\lP)^{2}$~\cite{Bousso1999}; the companion suspension paper adopts area counting for its
entropy hierarchy for exactly that reason~\cite{PB2026susp}.  The two counts answer different
questions---``how many field degrees of freedom fit at resolution $\varepsilon$'' versus ``how
much entropy can gravitate before collapse''---and their reconciliation (which degrees of
freedom are physical in a gravitating domain) is a genuine open problem of the framework, listed
in Sec.~\ref{sec:open}.  We flag the tension rather than hide it.
\end{remark}

\begin{proposition}[Causal integration bound]
\label{prop:coh}
On the compact causal spacetime $(\Sthree_{R}\times\R,\;g_{R}-c^{2}dt^{2})$, the causal
diameter is $D_{c}=\pi R$ and every event becomes causally accessible to every observer only
after the global accessibility time
\begin{equation}
t_{\mathrm{coh}}=\frac{\pi R}{c}
\end{equation}
\cite{PB2026light}.  Hence no observer can have integrated the full capacity of
Corollary~\ref{cor:planck-capacity} before $t_{\mathrm{coh}}$: capacity is a property of the
domain; \emph{integrated} information is bounded by causal accessibility.
\end{proposition}

\begin{definition}[Reconstructible information of a quasi-local domain]
\label{def:reconstructible}
The \emph{reconstructible information} of a quasi-local domain at time $T$ is the
distinguishability realized in the operational alphabet $\mathcal A_{\varepsilon}$, restricted
to the causally accessible region at $T$, above the grain of Postulate~\ref{post:grain}.  This
is a definition of the \emph{object}; given a source distribution on the alphabet, its
\emph{quantity} is measured by the standard Shannon functional.  The framework supplies the
alphabet; Shannon's theory then applies unchanged.
\end{definition}

\section{Dynamics of Carriers: Reconstruction Flows}
\label{sec:dynamics}

A theory of information must say not only how much a carrier holds but what happens to
distinguishability when the carrier itself is deformed.  The companion Ricci-type analysis
\cite{BP2026ricci} supplies the controlled laboratory:

\begin{enumerate}[label=(\alph*)]
\item \emph{Smoothing erases shape distinctions first.}  On the Berger family the flow drives
the anisotropy ratio $u\to1$: shape/fibre-squashing distinctions are the first casualties of
reconstruction.
\item \emph{Then the scale mode fails.}  After rounding, the unnormalized flow collapses the
remaining round mode; in reconstruction language, the variable has exhausted the distinguishable
structure of the domain.
\item \emph{The topological class survives where the metric fails.}  Throughout the smooth
regime the carrier remains $\Sthree$; $H_{3}(\Sthree)\simeq\Z$ is untouched by metric
degeneration until the description itself leaves the smooth category.  Metric erasure and
topological erasure are different operators with different prices
(Proposition~\ref{prop:erasure} versus curvature blow-up).
\item \emph{The grain closes the domain from below.}  The Planck barrier makes
$\{a\ge\lP\}$ forward-invariant with $|\mathrm{Rm}|^{2}\le C/\lP^{4}$: reconstruction never
assigns states below the grain, so the informational chain \eqref{eq:chain} never evaluates on
sub-grain configurations.
\end{enumerate}

\begin{remark}[Interpretation: two different limits of reconstructible information]
Backward, reconstruction meets the grain: excessive curvature, $a\in[\lP,\infty)$.  Forward, it
meets not contraction but \emph{insufficient distinguishability}: expansion, dilution, horizon
formation degrade the operational alphabet accessible to any observer
(Proposition~\ref{prop:coh} tightening against a growing domain).  Past limit: the $\lP$-boundary
of smooth geometry.  Future limit: loss of recoverable quasi-local structure.  Neither says
``reality ends''; each says this language of quasi-local informational reconstruction ceases to
have full meaning there~\cite{BP2026ricci}.
\end{remark}

\section{Relation to Existing Theories}
\label{sec:relation}

\paragraph{Shannon.}  Shannon's theory is untouched and presupposed downstream: it measures
sources on alphabets.  The present framework answers the prior question of where physical
alphabets come from (carrier + grain + causal access, Definition~\ref{def:reconstructible}) and
what operations on carriers preserve or destroy the distinctions the alphabet encodes.

\paragraph{Kolmogorov.}  Algorithmic information presupposes a symbol set and a machine; both
are carriers in the present sense, and Corollary~\ref{cor:no-invertible} applies to their
states.

\paragraph{Landauer--Bennett thermodynamics of computation.}  Anchor~\ref{anchor:landauer}: the
erasure discipline proved topologically in Sec.~\ref{sec:distinguishability} is realized
thermodynamically as the $k_{B}T\ln2$ bound.  The two statements are logically independent and
mutually reinforcing; neither is derived from the other here.

\paragraph{Quantum information.}  The qubit's phase--direction geometry is exactly the
$\Sthree\to\Stwo$ Hopf structure of the carrier chain (Sec.~\ref{sec:carriers}); the framework
locates quantum information's state space inside the same carrier hierarchy without re-deriving
quantum mechanics.

\paragraph{Wheeler.}  ``It from bit''~\cite{Wheeler1990} makes the bit primitive.  The present
framework inverts the emphasis without discarding the insight: \emph{no bit without carrier; no
carrier without distinguishability}---and neither pole is a self-standing substance
(Sec.~\ref{sec:interpretation}).  Information is not the ground of being here; it is the
invariant of a relational structure that becomes physical only through
\eqref{eq:chain}.

\section{Interpretive Layer}
\label{sec:interpretation}

Everything in this section is labelled \emph{Interpretation}: a reading discipline, not physics.

\begin{remark}[Non-emergence]
The theory does not assert the absolute emergence of information, space, or geometry.
$\Szero$ gives minimal distinguishability; suspension changes its carriers
(Theorem~\ref{thm:transport}); $\Sthree$ gives the special compact spin/Hopf carrier; the Planck
scale gives the lower metric-causal boundary of physical reconstruction
(Postulate~\ref{post:grain}).  Information is not produced from non-being and is not destroyed
inside admissible structural transformations; it changes its carrier of distinguishability.
Cessation of a conditioned representation is not annihilation of the invariant; it is release
from a conditioned carrier.
\end{remark}

\begin{remark}[Two protections]
First, against self-being: no phenomenon in the chain---distinction, carrier, alphabet,
observer---is posited as self-existing; each is given by a construction
(cf.\ the tetralemma discipline~\cite{Garfield1995}).  Second, against substantializing
dependence itself: having refuted self-being, one may not erect ``Interdependence'' as a new
substance.  The first part cuts off the self-being of the phenomenon; the second cuts off the
self-being of dependence itself.  In mathematical register: empty of own-nature, not empty of
dependent functioning.
\end{remark}

\begin{remark}[The axis]
\emph{The Planck boundary is not the beginning of the universe; it is the universal lower
boundary of quasi-local geometric reconstruction}~\cite{BP2026ricci}.  In informational terms:
not a first bit, but the grain below which the very notion of an operational bit is not
assigned.
\end{remark}

\begin{tcolorbox}[colback=gray!4,colframe=gray!50,title=Core principle (Interpretation)]
Information is dependently reconstructible distinguishability: an invariant of relational
structure (Definition~\ref{def:profile}), preserved under structure-preserving re-description
(Theorem~\ref{thm:invariance}), transported without loss across carriers
(Theorem~\ref{thm:transport}), destroyed only by exhibited non-invertible operators at a price
(Proposition~\ref{prop:erasure}, Anchor~\ref{anchor:landauer}), physically realized only
through carrier, grain, and causal order (Postulates~\ref{post:carrier}--\ref{post:grain}), and
operationally finite on compact domains (Theorem~\ref{thm:capacity}).  Not self-being, not
nothing: dependently reconstructible distinguishability passing through a chain of carriers.
\end{tcolorbox}

\section{Falsifiability and Observational Contact}
\label{sec:falsifiability}

The mathematical layer is proved, not testable.  The operational layer is testable at its
postulates and at the geometric infrastructure it borrows.

\begin{enumerate}[label=(\arabic*)]
\item \emph{Compact-topology searches.}  If the quasi-local domain admits compact
$\Sthree$-type reconstruction, the mode basis is \eqref{eq:capacity}'s hyperspherical basis;
matched-circle and correlation searches in the CMB constrain the admissible
radii~\cite{Cornish1998,Planck2018}.  A confirmed detection would fix the alphabet; continued
nulls push $R$ upward without falsifying the framework's mathematical core.
\item \emph{Log-periodic modulation.}  The golden-ratio log-periodic ansatz
$\omega=2\pi/\ln\golden$ associated with the $\golden$-hierarchy of the wider programme has been
directly tested by us against DESI~DR2 BAO and Pantheon$+$ data with null results
\cite{BP2026null}; accordingly, no statement in the present paper depends on that ansatz, and we
record its current status as unsupported.  A theory of information that inherits a disfavored
ansatz silently would violate its own erasure discipline.
\item \emph{Thermodynamic contact.}  Single-bit Landauer dissipation has been verified
experimentally~\cite{Berut2012}; this anchors the erasure discipline of
Sec.~\ref{sec:distinguishability} at the laboratory scale.
\item \emph{Grain postulate.}  Postulate~\ref{post:grain} is presently anchored geometrically
\cite{BP2026ricci}, not verified experimentally; any operational sub-Planck distinguishability
would falsify it and with it Corollary~\ref{cor:planck-capacity}.
\end{enumerate}

\section{Conclusion and Open Problems}
\label{sec:open}

The paper set out to define information before using it.  The result is a three-layer theory:
a proved mathematical core (invariance, attributed erasure, carrier transport), an operational
layer with exact capacity and accessibility formulas conditional on two explicit postulates, and
an interpretive layer strictly quarantined from both.  The contribution claimed is exactly
this: a definition-first, theorem-based, falsifiability-aware account of what physical
information is, where its alphabets come from, what preserves it, what destroys it and at what
price, and how much of it a compact quasi-local domain can carry and causally integrate.

\begin{enumerate}[label=\textbf{OP\arabic*.},leftmargin=2.2em]
\item \textbf{Categorical formulation.}  Formalize carriers and operations as a category with a
distinguished class of forgetting morphisms; the profile $\Dd$ as a functor; erasure prices as a
cost structure on non-invertible morphisms.
\item \textbf{Quantum profile.}  Replace components by orthogonality: develop the
distinguishability profile for state spaces (flag manifolds, $\CP^{n}$), with the Hopf class as
the obstruction datum, and compare with Holevo capacity.
\item \textbf{Capacity under reconstruction flow.}  Couple Theorem~\ref{thm:capacity} to the
Ricci-type flow of~\cite{BP2026ricci}: derive $dN_{\varepsilon}/d\tau$ along Berger and warped
trajectories and identify the informational signature of the Planck barrier.
\item \textbf{Volume--area reconciliation.}  Resolve the tension of
Remark~\ref{rem:holography}: determine which band-limited degrees of freedom remain physical in
a gravitating domain, and whether the operational count degrades to area scaling dynamically.
\item \textbf{Higher-degree erasure costs.}  Landauer's bound prices degree-$0$ erasure.  Is
there a thermodynamic price hierarchy for erasing degree-$k$ distinctions (cycles, enclosures),
e.g.\ in topological memories and topologically ordered media?
\item \textbf{Tensor and spinor alphabets.}  Extend \eqref{eq:capacity} to transverse-traceless
tensor and Dirac spectra on $\Sthree_{R}$ (the Dirac spectrum is
$\pm(n+3/2)/R$~\cite{PB2026susp}) and give the full graded capacity of the carrier.
\item \textbf{Integrated information dynamics.}  Combine Proposition~\ref{prop:coh} with
realistic source models to bound how fast any internal observer can integrate the domain's
capacity, and compare with horizon-entropy production rates.
\end{enumerate}

The theory may fail; but if it fails, it fails by calculation---at a postulate, at a capacity
formula, at an observational anchor---not by taste.

\begin{thebibliography}{99}

\bibitem{Shannon1948}
C.~E.~Shannon, \emph{A mathematical theory of communication},
Bell Syst.\ Tech.\ J.\ \textbf{27}, 379--423, 623--656 (1948).

\bibitem{Kolmogorov1965}
A.~N.~Kolmogorov, \emph{Three approaches to the quantitative definition of information},
Problems Inform.\ Transmission \textbf{1}, 1--7 (1965).

\bibitem{Landauer1961}
R.~Landauer, \emph{Irreversibility and heat generation in the computing process},
IBM J.\ Res.\ Dev.\ \textbf{5}, 183--191 (1961).

\bibitem{Landauer1991}
R.~Landauer, \emph{Information is physical}, Phys.\ Today \textbf{44}(5), 23--29 (1991).

\bibitem{Bennett1982}
C.~H.~Bennett, \emph{The thermodynamics of computation---a review},
Int.\ J.\ Theor.\ Phys.\ \textbf{21}, 905--940 (1982).

\bibitem{Berut2012}
A.~B\'erut, A.~Arakelyan, A.~Petrosyan, S.~Ciliberto, R.~Dillenschneider, and E.~Lutz,
\emph{Experimental verification of Landauer's principle linking information and thermodynamics},
Nature \textbf{483}, 187--189 (2012).

\bibitem{Wheeler1990}
J.~A.~Wheeler, \emph{Information, physics, quantum: the search for links}, in
\emph{Complexity, Entropy, and the Physics of Information}, ed.\ W.~H.~Zurek
(Addison-Wesley, 1990).

\bibitem{Hatcher2002}
A.~Hatcher, \emph{Algebraic Topology} (Cambridge University Press, 2002).

\bibitem{Bousso1999}
R.~Bousso, \emph{A covariant entropy conjecture}, JHEP \textbf{07}, 004 (1999).

\bibitem{Cornish1998}
N.~J.~Cornish, D.~N.~Spergel, and G.~D.~Starkman,
\emph{Circles in the sky: finding topology with the microwave background radiation},
Class.\ Quantum Grav.\ \textbf{15}, 2657--2670 (1998).

\bibitem{Planck2018}
Planck Collaboration, \emph{Planck 2018 results.\ VI.\ Cosmological parameters},
Astron.\ Astrophys.\ \textbf{641}, A6 (2020).

\bibitem{PB2026susp}
A.~Preece and B.~Batenin, \emph{Suspension of Spheres and the Minimal Role of $\Sthree$ in
Spin, Hopf Geometry, and Irreversible Hierarchy}, research manuscript (2026).

\bibitem{BP2026ricci}
B.~Batenin and A.~Preece, \emph{Hopf-Lifted Ricci-Type Flow on $\Sthree$: A No-Go for Bounded
and Engineered Forcing, and a Planck-Barrier Invariant Region}, research manuscript (2026).

\bibitem{PB2026light}
A.~Preece and B.~Batenin, \emph{Light as Causal Coherence Structure in Compact $\Sthree$
Geometry}, research manuscript (2026).

\bibitem{BP2026null}
B.~Batenin and A.~Preece, \emph{A Null Test of Golden-Ratio Log-Periodicity
($\omega=2\pi/\ln\varphi$) against DESI DR2 BAO and Pantheon$+$}, research manuscript (2026).

\bibitem{Garfield1995}
N\=ag\=arjuna, \emph{M\=ulamadhyamakak\=arik\=a}; English translation and commentary:
J.~L.~Garfield, \emph{The Fundamental Wisdom of the Middle Way} (Oxford University Press, 1995).

\end{thebibliography}

\end{document}